\section{데이터}
\label{sec:data}

\subsection{PUBG 경기 텔레메트리}
\label{sec:data:telemetry}

% P1: 무엇이 기록되는가
PUBG는 공개 API를 통해 경기마다 텔레메트리를 제공한다~\cite{PUBGAPI}.
텔레메트리는 경기에 참가한 모든 플레이어의 위치를 약 10초 간격으로 기록하며, 표본으로 추출한 약 1,738만 개의 연속 위치 간격에서 중앙값은 10초, 11초 이하의 비율은 99.93\%였다.
위치 기록과 함께 피해, 기절, 처치, 아이템의 획득·장착·사용, 안전구역의 상태 변화, 단계 전환, 착지와 같은 사건이 발생 시각과 함께 기록된다.
따라서 이 자료는 모든 팀의 이동 궤적과 팀 사이의 교전, 그리고 안전구역의 변화를 하나의 경기 안에서 함께 담는다.

% P2: 수집 규모와 희소성
본 연구는 PUBG 공개 API에서 수집한 스팀 서버의 경기 가운데 \ref{sec:data:benchmark}절의 기준으로 선정한 13,948경기를 분석에 사용한다.
이 경기들은 2026년 3월 26일부터 6월 2일까지 생성되었으며, 위치 표본 약 9,643만 건과 사건 약 2억 2,597만 건, 모두 약 3억 2,240만 건의 기록을 포함한다(표~\ref{tab:corpus}).
경기 종료 후의 통계를 집계한 공개 데이터셋은 존재하지만~\cite{KagglePUBG}, 경기 중 모든 플레이어의 위치를 이 해상도로 담은 자료는 드물며, 기존 배틀로얄 예측 연구가 사용한 경기 내 기록도 공개되지 않았다~\cite{Sagong25a,KRAFTON}.

\begin{table}[t]
\centering
\caption{분석에 사용한 13,948경기의 텔레메트리 기록 유형별 건수}
\label{tab:corpus}
\begin{tabular}{lr}
\hline
기록 유형 & 건수 \\
\hline
위치 표본 (약 10초 간격) & 96,426,398 \\
피해 & 90,178,937 \\
아이템 획득·버림 & 67,865,634 \\
아이템 장착 & 46,218,920 \\
아이템 사용 & 15,229,394 \\
안전구역 상태 갱신 & 2,134,448 \\
착지 & 1,673,902 \\
처치 & 1,267,179 \\
기절 & 1,199,290 \\
단계 전환 & 200,979 \\
\hline
합계 & 322,395,081 \\
\hline
\end{tabular}
% DB 확인: 벤치마크 match_id 13,948개로 COUNT (2026-10-06).
\end{table}

\subsection{학습 표본 구성}
\label{sec:data:samples}

% P3: 공통 시간축, 단계 경계, 인과적 절단
원시 기록을 학습 표본으로 바꾸려면 세 가지 시간 처리가 필요하다.
첫째, 모든 기록을 하나의 시간축에 놓아야 한다.
위치 기록에 함께 저장된 경기 경과 시간은 사건 기록의 타임스탬프와 기준시각이 달라 경기별로 최대 41.9초까지 어긋나므로, 본 연구는 모든 기록의 타임스탬프를 경기 시작 사건으로부터의 경과 시간으로 변환하여 하나의 시간축으로 사용한다.
둘째, 단계 경계는 단계 전환 사건으로 정한다.
텔레메트리는 새 안전구역의 공개와 구역 축소의 시작을 단계 전환 사건으로 기록하며, 한 구역으로의 축소가 끝나는 시각에 다음 구역이 공개된다~\cite{PUBGAPI}.
본 연구는 $k$번째 안전구역이 공개된 시점부터 그 구역으로의 축소가 끝날 때까지를 단계 $k$로 정의하고, 축소가 끝나 다음 안전구역이 공개되는 시각을 단계 $k$의 종료 시각 $b_k$로 둔다.
셋째, 예측 시점은 $b_k$로 두고, 모든 입력은 그보다 10초 앞선 특징 절단 시각 $a_k=b_k-10$초 이전의 기록만으로 구성한다.
이 제한은 특징뿐 아니라 생존 인원, 팀 사이의 관계, attention 마스크에도 적용되며, 생존 여부 역시 최종 탈락 시각이 아니라 $a_k$ 이전의 처치와 위치 기록으로 판정한다.
처치된 플레이어는 처치 후 20초 이상 지난 시점에 체력이 남은 위치 기록이 나타날 때에만 재투입된 것으로 본다.
12개 코호트에서 17경기씩 고른 204경기의 1,293개 단계 경계에서 $a_k$ 이후의 원시 기록을 모두 지우고 입력을 다시 구성한 결과, 모든 입력이 원래 입력과 일치하였다.
특징 절단 시각에 생존한 플레이어의 99.4\%는 직전 15초 안에 위치 표본을 가진다.

% P4: 사건 기반 탈락 라벨
탈락 라벨은 사건 기록으로 정의한다.
기절은 이후 소생이 가능한 상태이므로 최종 처치와 구분하며, 팀의 탈락 시각은 재투입되지 않은 마지막 팀원이 처치된 시각으로 정의한다.
이렇게 정한 탈락 시각은 모든 탈락 팀에서 게임이 팀의 최종 순위를 기록한 처치 사건의 시각과 일치하였다.
우승 팀은 탈락 사건 없이 경기가 끝나므로 경기 종료 시각에서 관측을 종료한다.
단계 $k$의 위험집합은 예측 시점 $b_k$에 생존한 팀들이며, 각 팀의 라벨은 다음 단계의 종료 시각 $b_{k+1}$까지 탈락하면 1, 그렇지 않으면 0이다.
다음 단계가 경기의 마지막 단계이면 $b_{k+1}$ 대신 경기 종료 시각을 사용한다.
경기 종료 기록이 없거나 라벨 구간에서 위치 기록이 60초 넘게 끊긴 경기는 없었다.
예측 시점은 단계 1부터 7까지의 경계에 둔다.
따라서 분석의 단위는 경기, 팀, 단계 경계의 조합인 팀--단계 표본이다.

\subsection{벤치마크}
\label{sec:data:benchmark}

% P5: 모집단 정제와 코호트
분석 대상은 일반 경기와 랭크 경기로 한정하였다.
실제 경쟁이 이루어지지 않는 튜토리얼과 훈련장 경기는 제외하였고, 일반적인 경기 양상과 다른 두 모집단도 제외하였다.
계정의 약 92\%가 AI인 airoyale 경기와, 경기마다 약 30개의 비플레이어 전투 계정이 기록되는 Vikendi 맵의 경기이다.
일반 경기에도 AI 계정이 섞여 있으므로, 각 팀의 AI 계정 비율을 입력 특징에 포함하였다.
경기는 맵, 시점(3인칭·1인칭), 팀 규모(1인·2인·4인)의 조합으로 코호트를 나누었으며, airoyale 경기를 제외하기 전 기준으로 300경기 이상인 조합을 코호트로 정하였다.
최종 벤치마크는 12개 코호트의 13,948경기이며, 이 중 11,285경기는 일반 경기, 2,663경기는 랭크 경기이다.
% 보류: 랭크 경기와 대회의 안전구역 설정 일치는 규정 원문 또는 경기 설정 기록으로 확인한 뒤 복원.
% 랭크 경기는 e스포츠 대회와 같은 안전구역 설정을 사용한다~\cite{PUBGRules}.

% P6: 분할과 표본 수
각 코호트의 경기는 경기 단위로 학습, 검증, 시험 집합에 8:1:1로 나누었으며, 분할은 고정하여 모든 비교에 같은 분할을 사용한다.
그 결과 학습, 검증, 시험 집합은 각각 11,188경기, 1,376경기, 1,384경기이며, 코호트의 크기는 224경기부터 4,159경기까지이다.
위험집합에 속한 팀--단계 표본은 모두 817,766개이고, 그중 217,804개(26.6\%)가 다음 단계 안에 탈락하였다.
다음 단계 탈락 비율은 단계 1의 19.6\%에서 단계 7의 50.4\%로 높아진다.
